\documentclass[a4paper,12pt]{article}

\usepackage[T1]{fontenc}
\usepackage[utf8]{inputenc}
\usepackage[brazil]{babel}
\usepackage[a4paper,margin=2.5cm]{geometry}

\begin{document}

\title{Capacitação: Construindo um Encurtador de URLs com Node.js e Express}
\author{}
\date{}

\maketitle

\section{Estrutura da Capacitação (1 hora)}

\subsection{Parte 1 --- Introdução (5 min)}

\subsubsection{O problema}

Pergunte:

\begin{quote}
Como o Bitly transforma uma URL enorme em algo pequeno?
\end{quote}

Exemplo:

\begin{verbatim}
https://www.youtube.com/watch?v=ABC123XYZ456
\end{verbatim}

vira

\begin{verbatim}
https://bit.ly/abc123
\end{verbatim}

Explique que o sistema:

\begin{itemize}
    \item Recebe uma URL;
    \item Gera um código;
    \item Salva a relação código $\rightarrow$ URL;
    \item Redireciona quando alguém acessa o código.
\end{itemize}

\subsection{Parte 2 --- O que é uma API REST? (5 min)}

Mostre:

\begin{verbatim}
Cliente
   |
   v
Requisição HTTP
   |
   v
Servidor
   |
   v
Resposta HTTP
\end{verbatim}

Exemplos:

\begin{verbatim}
GET /users
\end{verbatim}

Buscar dados.

\begin{verbatim}
POST /users
\end{verbatim}

Criar dados.

Explique rapidamente:

\begin{center}
\begin{tabular}{|c|c|}
\hline
Método & Função \\
\hline
GET & Buscar \\
POST & Criar \\
PUT & Atualizar \\
DELETE & Remover \\
\hline
\end{tabular}
\end{center}

\subsection{Parte 3 --- Criando o Projeto (10 min)}

\subsubsection{Inicialização}

\begin{verbatim}
npm init -y
\end{verbatim}

\subsubsection{Instalação}

\begin{verbatim}
npm install express
\end{verbatim}

Mostre o que apareceu:

\begin{itemize}
    \item node\_modules
    \item package.json
    \item package-lock.json
\end{itemize}

\subsection{Parte 4 --- Criando o Servidor (10 min)}

Arquivo:

\begin{verbatim}
const express = require("express");
const app = express();

app.listen(3000, () => {
    console.log("Servidor rodando");
});
\end{verbatim}

Explique:

\begin{verbatim}
express()
\end{verbatim}

Cria a aplicação.

Porta:

\begin{verbatim}
app.listen(3000)
\end{verbatim}

Pergunte:

\begin{quote}
O navegador conversa com quem?
\end{quote}

Resposta:

\begin{verbatim}
localhost:3000
\end{verbatim}

\subsection{Parte 5 --- Primeira Rota (5 min)}

\begin{verbatim}
app.get("/", (req, res) => {
    res.send("Olá mundo");
});
\end{verbatim}

Explique:

\begin{description}
    \item[req] Requisição.
    \item[res] Resposta.
\end{description}

Teste no navegador:

\begin{verbatim}
http://localhost:3000
\end{verbatim}

\subsection{Parte 6 --- Organizando Rotas (5 min)}

Mostre o arquivo:

\begin{verbatim}
routes/Users.js
\end{verbatim}

\begin{verbatim}
const router = express.Router();
\end{verbatim}

Explique:

Antes:

\begin{verbatim}
app.get(...)
app.post(...)
app.get(...)
\end{verbatim}

Tudo no mesmo arquivo.

Depois:

\begin{verbatim}
router.get(...)
router.post(...)
\end{verbatim}

Mais organizado.

\subsection{Parte 7 --- Criando o Encurtador (15 min)}

\subsubsection{Rota POST}

\begin{verbatim}
router.post("/shorten", ...)
\end{verbatim}

Recebe:

\begin{verbatim}
{
  "url": "https://google.com"
}
\end{verbatim}

\subsubsection{Middleware}

\begin{verbatim}
app.use(express.json());
\end{verbatim}

Explique:

Sem isso:

\begin{verbatim}
req.body
\end{verbatim}

vem vazio.

\subsubsection{Validação}

\begin{verbatim}
if (!url)
\end{verbatim}

Explique a importância de validar entradas.

\subsubsection{Gerando o Código}

\begin{verbatim}
console.log(
    Math.random()
        .toString(36)
        .substring(2, 8)
);
\end{verbatim}

Saídas possíveis:

\begin{verbatim}
ab12cd
xk92lm
qwerty
\end{verbatim}

\subsubsection{Banco em Memória}

\begin{verbatim}
const database = {};
\end{verbatim}

Exemplo:

\begin{verbatim}
database["abc123"] = "https://google.com";
\end{verbatim}

Resultado:

\begin{verbatim}
{
   abc123: "https://google.com"
}
\end{verbatim}

\subsubsection{Retorno}

\begin{verbatim}
res.json({
    shortUrl: "http://localhost:3000/abc123"
});
\end{verbatim}

Mostre no Postman ou Insomnia.

\subsection{Parte 8 --- Redirecionamento (5 min)}

\begin{verbatim}
router.get("/:code")
\end{verbatim}

Quando acessar:

\begin{verbatim}
localhost:3000/abc123
\end{verbatim}

Busca:

\begin{verbatim}
database["abc123"]
\end{verbatim}

Se encontrar:

\begin{verbatim}
res.redirect(originalUrl);
\end{verbatim}

Explique o código HTTP 302.

\section{Demonstração Final}

\subsection{1. Criar URL}

\begin{verbatim}
POST /shorten
\end{verbatim}

Body:

\begin{verbatim}
{
  "url": "https://www.google.com"
}
\end{verbatim}

Resposta:

\begin{verbatim}
{
  "shortUrl": "http://localhost:3000/abc123"
}
\end{verbatim}

\subsection{2. Acessar}

\begin{verbatim}
http://localhost:3000/abc123
\end{verbatim}

Resultado:

\begin{verbatim}
Google
\end{verbatim}

\end{document}